% Generated draft: requires mathematical and source review.
\documentclass[11pt]{article}
\usepackage[margin=25mm]{geometry}
\usepackage{fontspec}
\setmainfont{Latin Modern Roman}
\usepackage{amsmath,amssymb,mathtools,unicode-math}
\setmathfont{Latin Modern Math}
\usepackage{enumitem,needspace,xurl,hyperref}
\setlength{\parindent}{0pt}
\setlength{\parskip}{5pt}
\setlength{\emergencystretch}{3em}
\newcommand{\N}{\mathbb{N}}
\newcommand{\Z}{\mathbb{Z}}
\newcommand{\Q}{\mathbb{Q}}
\newcommand{\R}{\mathbb{R}}
\newcommand{\C}{\mathbb{C}}
\newcommand{\F}{\mathbb{F}}
\newcommand{\A}{\mathbb{A}}
\newcommand{\PP}{\mathbb{P}}
\newcommand{\E}{\mathbb{E}}
\newcommand{\eps}{\varepsilon}
\newcommand{\dd}{\,\mathrm{d}}
\begin{document}
\section*{UnsolvedMath ID: 30002270}
\textbf{Source problem code:} catalogue entry
\subsection{Definitions and mathematical statement}
For a finite simple undirected graph \(G=(V,E)\), let \(d_G(v)=|\{u\in V:uv\in E\}|\), let \(\chi(G)\) be its chromatic number, let \(G[S]\) be the subgraph induced by \(S\subseteq V\), and let \(\omega(H)\) be the maximum size of a clique in \(H\). A conventional interpretation of the source notation is the closed neighborhood
\[
\widetilde N(v)=\{v\}\cup\{u\in V:uv\in E\}.
\]
The supplied record does not explicitly specify the finite simple undirected graph setting or define \(\widetilde N(v)\); these conventions therefore require confirmation.

With this interpretation, define
\[
\gamma_\ell(v)=\frac12\bigl(d_G(v)+1+\omega(G[\widetilde N(v)])\bigr).
\]
The supplied conjectural assertion is that every graph \(G\) satisfies both
\[
\chi(G)\leq\max_{v\in V}\frac12\bigl(d_G(v)+1+\omega(G[\widetilde N(v)])\bigr)
\]
and
\[
\chi(G)\leq\max_{uv\in E}\frac12\bigl(\gamma_\ell(u)+\gamma_\ell(v)\bigr).
\]
These displays reproduce the unrounded bounds in the record. The accompanying literature assessment calls the target an integer bound, so whether ceiling operations are missing requires source review; no rounding has been inserted here. Under the finite graph interpretation, both maxima are defined when \(E\ne\varnothing\). The source's universal quantifier also includes graphs with no edges unless its domain is restricted, but no convention for the maximum over an empty edge set is supplied. For the graph with no vertices, the vertex maximum is also unspecified. Consequently, the universal assertion remains unresolved as a precise formulation.
\subsection{Short English statement}
The record proposes local and edgewise upper bounds on the chromatic number, using degrees and clique numbers in neighborhoods. The displayed formulas are preserved exactly in their unrounded form. The graph domain, neighborhood convention, possible missing ceilings, and treatment of graphs with no edges require review. The record labels the problem open; this is supplied curation rather than an independently verified status.
\subsection{Sources}
\begin{itemize}
\item Local UnsolvedMath registry, permanent ID \(30002270\), title \emph{Local and Edgewise Reed Bounds}, source problem code \texttt{catalogue entry}, in the Oberwolfach Reports Open Problems collection. The record gives proposed year 2013 and identifies no proposer.
\item Supplied report citation: \emph{Graph Theory} (2013), Oberwolfach Reports, pp.~67--128, DOI \texttt{10.4171/owr/2013/02}. \url{https://doi.org/10.4171/owr/2013/02}.
\item Supplied literature reference: Edwards--King, \emph{A superlocal version of Reed's conjecture}. \url{https://doi.org/10.37236/2666}.
\item Literature assessment dated exactly \texttt{2026-08-21}: the supplied curation labels the problem open, identifies the superlocal integer bound as Conjecture~3, and reports fractional and special-class results but no general proof found. This is supplied dated curation, not independently verified. The record also labels its statement \texttt{corrected\_verified} and reports checking against the official Oberwolfach report; that verification claim is supplied metadata and has not been independently checked here.
\end{itemize}
\end{document}
